\section{Task and Evaluation Specification}
\label{app:task_details}

\subsection{Piecewise-stationary modes}

Table~\ref{tab:task_modes} specifies the v5 MO-Gymnasium tasks. Each mode
rescales the base MuJoCo parameters; unlisted multipliers are one and reward
shifts are zero. For base reward $r_j^{\mathrm{base}}$, scale $\rho_j$ and a
finite nonzero cap parameter $c_j$, the transformed coordinate is
\begin{equation}
 r_j=|c_j|\tanh\!\left(
      \frac{\rho_j r_j^{\mathrm{base}}}{|c_j|}\right),
 \label{eq:mode_reward_transform}
\end{equation}
An infinite cap leaves $\rho_j r_j^{\mathrm{base}}$ unchanged. Finite caps
give smooth, symmetric saturation with the tabulated magnitudes $|c_j|$.

\begin{table}[htbp]
\centering\small
\caption{Mode parameters used by both methods. $g$, $m$, $f$ and $d$ multiply
gravity, body masses, geom friction and joint damping, respectively.
Reward coordinates follow the base task order.}
\label{tab:task_modes}
\begin{tabular}{@{}llccc@{}}
\toprule
Task & Mode & Physics changes & Reward scale & Soft-cap magnitudes\\
\midrule
\multirow{4}{*}{HalfCheetah}
 & Normal & --- & $(1,1)$ & $(\infty,\infty)$\\
 & Low speed & $m=2,\ g=0.6$ & $(1,1)$ & $(0.3,\infty)$\\
 & Low cost & $m=0.5,\ g=1.5$ & $(1,1)$ & $(\infty,0.05)$\\
 & Locked & $d=3$ & $(1,1)$ & $(1,0.2)$\\
\midrule
\multirow{4}{*}{Hopper}
 & Normal & --- & $(1,1,1)$ & $(\infty,\infty,\infty)$\\
 & Fragile & $m=1.5$ & $(1,1,1)$ & $(\infty,0.5,\infty)$\\
 & Cheap energy & --- & $(1,1,1)$ & $(\infty,\infty,0.05)$\\
 & Grounded & $g=1.5$ & $(1,1,1)$ & $(0.5,0.3,\infty)$\\
\midrule
\multirow{3}{*}{Ant}
 & Energy save & $m=2,\ f=1.5$ & $(0.1,0.1,5)$ & $(\infty,\infty,\infty)$\\
 & Race & $g=0.7,\ f=0.5$ & $(0.1,5,0.1)$ & $(\infty,\infty,\infty)$\\
 & Cautious & $g=1.5,\ f=0.3$ & $(5,0.1,0.2)$ & $(\infty,\infty,\infty)$\\
\bottomrule
\end{tabular}
\end{table}

Coordinates are forward velocity and negative control cost on HalfCheetah;
forward velocity, jumping height and negative control cost on Hopper; and
$x$ velocity, $y$ velocity and negative control cost on Ant. Hopper's base
survival bonus and Ant's base survival/contact terms enter every coordinate
before the mode transformation.

Training starts in mode zero. Dwell lengths are
$D=\max\{1000,\lfloor X\rfloor\}$ for $X$ exponentially distributed with
scale 25,000; the next mode is uniform among the other modes. Episode resets
preserve the mode and switching clock. The independently seeded schedule gives
paired methods identical mode histories. The controller observes states and
rewards, not mode labels or historical anchor metadata.

\subsection{Locked-mode final-policy evaluation}

Evaluation uses deterministic actions and, for BAMOR, fresh context trackers,
fixing the mode and requested preference within each episode.
The two-objective grid is
$W_2=\{(i/20,1-i/20):i=0,\ldots,20\}$; the three-objective grid is
$W_3=\{(i/5,j/5,(5-i-j)/5):i,j\in\mathbb Z_{\geq0},\ i+j\leq5\}$.
Both contain 21 weights, including vertices. Every weight in mode $h$ uses
the three reset seeds $100000+1000s+10h+i$, with training seed $s$ and
$i\in\{0,1,2\}$. Return vectors are averaged over these episodes before
computing direct utility or HV. Termination or a 1,000-step limit ends an
episode. Evaluation performs no parameter updates and contributes no samples
to training replay or operator scores.

The fixed HV reference is $(-50,-3000)$ on HalfCheetah and
$(-50,-100,-3000)$ on Hopper and Ant. HV is computed within each mode and
then averaged across modes. At 75k each method evaluates one final actor; at
50k each evaluates 19 retained actors. BAMOR pairs all its actors with the
final encoder. Archive evaluation describes a policy set, not an online
learning curve.
